\documentclass{article}
\usepackage{amsmath, amssymb, amsthm}
\usepackage{hyperref}
\usepackage{geometry}
\geometry{margin=1in}

\title{Erd\H{o}s Problem \#327}
\date{Last updated: 2025-08-31}
\author{Source: erdosproblems.com}

\begin{document}
\maketitle

\noindent\textbf{Status:} OPEN \quad \textbf{No prize}

\noindent\textbf{Tags:} number theory, unit fractions

\medskip

\noindent\textbf{Problem Statement:}

Suppose $A\subseteq \{1,\ldots,N\}$ is such that if $a,b\in A$ and $a\neq b$ then $a+b\nmid ab$. Can $A$ be 'substantially more' than the odd numbers? What if $a,b\in A$ with $a\neq b$ implies $a+b\nmid 2ab$? Must $\lvert A\rvert=o(N)$?




The connection to unit fractions comes from the observation that $\frac{1}{a}+\frac{1}{b}$ is a unit fraction if and only if $a+b\mid ab$. Wouter van Doorn has given an elementary argument that proves that if $A\subseteq \{1,\ldots,N\}$ has $\lvert A\rvert \geq (25/28+o(1))N$ then $A$ must contain $a\neq b$ with $a+b\mid ab$ (see the discussion in [301] ). See also [302] .

\noindent\textbf{Related OEIS sequences:} A384927


\bigskip
\noindent\small{Source: \url{https://www.erdosproblems.com/327}}
\end{document}
